\chapter{Strumenti, ambiente e codice}
\label{app:strumentiecodice}

\section{Strumenti e ambiente di simulazione - trattazione estesa}
\label{app:strumenti}
\begin{onehalfspace}

La scelta degli strumenti di simulazione non è neutrale. In informatica quantistica, il framework adottato determina quali modelli di rumore sono disponibili, verso quale hardware è possibile traslare i circuiti e con quali risultati della letteratura è possibile confrontarsi direttamente. Questa sezione descrive gli strumenti utilizzati nella parte sperimentale di questa tesi, motivando ciascuna scelta e collocandola rispetto alle alternative disponibili.

% -----------------------------------------------
\subsection{Qiskit: il Framework Quantistico di Riferimento}
\label{sec:qiskit}

\textbf{Qiskit} (\textit{Quantum Information Science Kit}) è un \ac{SDK} open-source sviluppato da \ac{IBM} Quantum per la programmazione, la simulazione e l'esecuzione di algoritmi quantistici \cite{qiskit2023}. Rilasciato pubblicamente nel 2017, offre integrazione diretta con l'hardware \ac{IBM} ed è uno dei framework ampiamente usati nella ricerca su circuiti quantistici a corto termine \cite{skosana2021, shor_hpec2023}.

\subsubsection{Architettura del Framework}

Qiskit è organizzato in tre componenti principali, ciascuno con un ruolo distinto nel ciclo di vita di un esperimento quantistico:

\begin{itemize}

    \item \textbf{Qiskit (core)} — il nucleo del framework, già noto come Qiskit Terra. Fornisce le primitive per costruire i circuiti, compilarli verso un backend target tramite \texttt{transpile} e accedere alle interfacce \texttt{Sampler} ed \texttt{Estimator}. Gestisce la pipeline dalla descrizione astratta fino all'invio al simulatore o all'hardware reale.

    \item \textbf{Qiskit Aer} — il modulo di simulazione classica ad alte prestazioni. Implementa diversi metodi di simulazione esatta e approssimata, con supporto nativo per modelli di rumore configurabili. È il componente centrale di questa tesi e viene descritto in dettaglio nella Sezione~\ref{sec:qiskit_aer}.

    \item \textbf{Qiskit IBM Runtime} — l'interfaccia per l'esecuzione su processori \ac{IBM} Quantum reali tramite cloud (\texttt{qiskit-ibm-runtime}). Fornisce le primitive \texttt{Sampler} e \texttt{Estimator} ottimizzate per l'hardware reale, con supporto per tecniche di soppressione del rumore come il \textit{Dynamical Decoupling} e il \textit{Gate Twirling} \cite{ibm_shor_tutorial}.

\end{itemize}

\subsubsection{Perché Qiskit e non un'alternativa}
\label{sec:confronto_framework}

La scelta di Qiskit, maturata anche dall'analisi del tutorial ufficiale \ac{IBM} per l'algoritmo di Shor~\cite{ibm_shor_tutorial}, è motivata da quattro ragioni convergenti.

\textbf{Prima ragione: controllo esplicito del modello.} Qiskit Aer permette di comporre canali depolarizzanti, termici e di readout e di legarli a un gate set dichiarato. Questa flessibilità consente esperimenti controllati e riproducibili; non rende però i preset uniformi della tesi direttamente equivalenti alla calibrazione di una QPU reale.

\textbf{Seconda ragione: copertura delle primitive necessarie.} Qiskit Aer espone nello stesso ambiente i canali depolarizzanti, termici, di readout e coerenti richiesti dal disegno. Altri framework, fra cui Cirq \cite{cirq2024} e PennyLane \cite{pennylane2022}, offrono anch'essi simulazione rumorosa; la scelta non pretende di stabilire una graduatoria generale, ma riduce le conversioni rispetto al codice Qiskit adottato.

\textbf{Terza ragione: disponibilità dell'implementazione di riferimento.} Il circuito dell'algoritmo di Shor utilizzato in questa tesi è tratto da un'implementazione open-source in Qiskit \cite{github_shor}. Questo punto di partenza, validato dalla letteratura \cite{skosana2021}, ha reso naturale mantenere Qiskit come framework unico lungo l'intero pipeline sperimentale.

\textbf{Quarta ragione: allineamento con le fonti implementative.} Diversi lavori usati come riferimento adottano Qiskit \cite{skosana2021, shor_hpec2023, ml_mitigation2023}. Condividere il framework facilita il confronto delle API e dei circuiti, ma non rende direttamente comparabili risultati ottenuti con gate set, versioni, transpiler o noise model differenti.

\begin{table}[H]
\centering

\small
\begin{tabularx}{\textwidth}{@{}>{\raggedright\arraybackslash}p{2.5cm}XXX@{}}
\toprule
\textbf{Framework} & \textbf{Sviluppatore} & \textbf{Simulatore con rumore} & \textbf{Target hardware primario} \\
\midrule
Qiskit + Aer        & IBM Quantum           & Completo (T1, T2, depolarz., readout) & IBM superconduttori \\
\addlinespace[0.3em]
Cirq                & Google Quantum AI     & Parziale                               & Google superconduttori \\
\addlinespace[0.3em]
PennyLane           & Xanadu                & Limitato                               & Vari (plug-in) \\
\addlinespace[0.3em]
Forest / PyQuil     & Rigetti               & Parziale                               & Rigetti superconduttori \\
\bottomrule
\end{tabularx}
\caption{Confronto tra i principali framework per la simulazione quantistica con rumore.}\label{tab:framework_confronto}

\end{table}

% -----------------------------------------------
\subsection{Qiskit Aer: il Simulatore}
\label{sec:qiskit_aer}

Qiskit Aer è il modulo di simulazione classica di Qiskit. Il suo componente principale è \texttt{AerSimulator}, un simulatore ad alte prestazioni che supporta più metodi di simulazione selezionabili a seconda della dimensione del circuito e del tipo di analisi richiesta.

\subsubsection{Metodi di Simulazione}

\texttt{AerSimulator} espone i seguenti metodi principali tramite il parametro \texttt{method}:

\begin{itemize}
    \item \texttt{statevector} — simulazione del vettore di stato completo. Richiede $2^n$ ampiezze complesse in memoria ed è il metodo usato dalla campagna rumorosa principale su $N=15$.

    \item \texttt{density\_matrix} — evolve la matrice densità $\rho$ anziché il vettore di stato, permettendo di modellare canali di rumore non unitari in modo esatto. Richiede $4^n$ elementi e riduce il limite pratico a circa 25 qubit.

    \item \texttt{stabilizer} — simulazione efficiente di circuiti Clifford puri (senza porte non-Clifford come $T$ o $R_z$ arbitrarie). Scala a migliaia di qubit ma non è applicabile all'algoritmo di Shor, che utilizza rotazioni di fase arbitrarie.

    \item \texttt{matrix\_product\_state} — rappresenta lo stato come rete di tensori; senza troncamenti imposti resta una rappresentazione controllata, ma il costo cresce con l'entanglement. È usato in M8/M13 per contenere il costo delle molte repliche.
\end{itemize}

La scelta è quindi per campagna, non globale: \texttt{statevector} per la prima fase N15 e \texttt{matrix\_product\_state} per la sensibilità M8/M13. Le istanze N21/N35 sono validate idealmente ma escluse dalle campagne rumorose generali.

\subsubsection{Modelli di Rumore}

Il modulo \texttt{qiskit\_aer.noise} fornisce le primitive per costruire modelli di rumore compositi. Ogni errore è rappresentato come un canale quantistico che viene applicato ai gate specificati, con la possibilità di assegnare errori diversi a qubit diversi (rumore eterogeneo) o un errore uniforme a tutti i qubit (\textit{all-qubit error}). I tipi di errore disponibili includono:

\begin{itemize}
    \item \texttt{depolarizing\_error} — errore depolarizzante su gate a 1 o 2 qubit, parametrizzato da $\lambda$ nella convenzione Aer; la probabilità Pauli non-identità è $3\lambda/4$ o $15\lambda/16$.
    \item \texttt{thermal\_relaxation\_error} — modella il rilassamento verso lo stato fondamentale ($T_1$) e la perdita di coerenza di fase ($T_2$), con un tempo di gate configurabile.
    \item \texttt{ReadoutError} — matrice di confusione che descrive la probabilità di misurare $\ket{0}$ come $\ket{1}$ e viceversa.
    \item Errori coerenti — applicazione di una rotazione indesiderata su ogni gate, modellabile tramite operatori unitari arbitrari.
\end{itemize}

Il modello di rumore usato nella prima fase combina errori depolarizzanti, rilassamento termico e readout error con parametri uniformi. È un modello illustrativo per isolare i meccanismi, non uno snapshot di \texttt{ibm\_marrakesh}: non include coupling map, routing, errori per qubit, crosstalk o deriva temporale. La configurazione dettagliata è riportata nel Capitolo~\ref{chap:obiettivi}.

\subsubsection{Transpilazione}

Prima dell'esecuzione, la funzione \texttt{transpile} riscrive il circuito nella base esplicita \texttt{rz/sx/x/cx} e applica ottimizzazioni strutturali. La campagna usa il livello di ottimizzazione 2 e il seed del transpiler 20260819; questa coppia, insieme all'hash della sequenza compilata, è registrata in ogni artefatto. La base esplicita impedisce che porte composte restino opache e sfuggano al canale di rumore previsto.

% -----------------------------------------------
\subsection{Accelerazione GPU con Qiskit Aer}
\label{sec:gpu}

L'accelerazione tramite \ac{GPU} NVIDIA (libreria \texttt{qiskit-aer-gpu}, scheda
GeForce RTX 4070 da 12 GB) era stata pianificata per la generazione dei dataset e per
le campagne replicate. Su \ac{WSL} il driver CUDA riconosceva la scheda, ma il backend
di Aer sollevava a runtime l'errore \texttt{Simulation device "GPU" is not supported on
this system}. \textbf{Tutti gli esperimenti riportati nella tesi sono quindi stati
eseguiti su \ac{CPU}}, e il backend effettivo è registrato nei metadati di ogni
artefatto (Sezione~\ref{sec:gpu_limit}).

% -----------------------------------------------
\subsection{scikit-learn: la Libreria per il Classificatore}
\label{sec:scikit}

Il Metodo 2 richiede l'addestramento di un classificatore binario. La libreria scelta è scikit-learn \cite{scikit_learn}, una delle librerie di machine learning più consolidate nell'ecosistema Python scientifico.

La motivazione principale è pragmatica: scikit-learn espone un'interfaccia unificata (\texttt{fit} / \texttt{predict}) per un ampio catalogo di classificatori — Random Forest, \ac{SVM}, \ac{MLP}, e altri — permettendo di confrontarli sullo stesso dataset con minime modifiche al codice. Questo facilita la selezione sperimentale del modello migliore senza dover introdurre dipendenze da framework più pesanti (TensorFlow, PyTorch) per un classificatore che, per il problema in esame, si prevede non richieda architetture profonde.

% -----------------------------------------------
\subsection{Ambiente di Esecuzione: WSL con Ubuntu}
\label{sec:wsl}

Le simulazioni vengono eseguite in ambiente Linux tramite \ac{WSL} (versione 2) su Windows 11. La scelta è motivata dalla necessità di mantenere la compatibilità con l'ecosistema Linux di Qiskit e delle librerie scientifiche Python, garantendo al contempo l'integrazione con i driver CUDA per \ac{WSL} forniti da NVIDIA.

\ac{WSL} 2 esegue un kernel Linux completo in una macchina virtuale leggera.

\subsubsection{Configurazione riproducibile e versioni}
\label{app:ambiente_versioni}

La campagna canonica usa Windows~11 come sistema ospite, \ac{WSL}2 con Ubuntu~24.04,
Python~3.12 e Qiskit~2.5; Stim~1.16.0 e PyMatching~2.4.0 sono impiegati nella parte
surface-code. L'ambiente virtuale è creato e popolato dal file
\texttt{requirements.txt} come nel
Listato~\ref{lst:installazione_dellambiente_python_e}. Ogni artefatto registra inoltre
le versioni effettivamente importate di Qiskit, Aer, NumPy, SciPy e scikit-learn, così
che il file dei requisiti e il manifest dell'esecuzione svolgano ruoli complementari:
il primo prescrive l'ambiente, il secondo certifica quello realmente usato. Tutte le
campagne riportate sono state orchestrate su CPU.

\end{onehalfspace}

% -----------------------------------------------

\section{Codice sorgente dei componenti implementati}
\label{app:codice}
% Listati della pipeline sperimentale. I sorgenti eseguibili restano l'autorità;
% gli estratti mostrano il contratto versione 2 senza incorporare risultati.

\begin{onehalfspace}

Questa sezione documenta il contratto di esecuzione, l'organizzazione dei moduli, i
dettagli dei codici QEC e gli estratti dei sorgenti. Gli output numerici non sono codificati
nei listati: ogni campagna li salva in un artefatto JSON schema~2 corredato di manifest.

\subsection{Contratto implementativo e riproducibilità}
\label{app:contratto_codice}

\subsubsection{Netlist compilata e manifest}
\label{app:netlist_manifest}

La compilazione è centralizzata e condivisa da training, baseline e sweep. Il contratto
fissa \texttt{optimization\_level=2}, seed del transpiler e base nativa
\texttt{rz/sx/x/cx}. Con Qiskit~2.5, $N=15$, $a=7$ e otto qubit di conteggio, la
netlist ha profondità 412 e contiene 224 porte \texttt{cx}; il costo molto maggiore
delle istanze $N=21$ e $N=35$ è riportato nella Sezione~\ref{app:fase1_contratto}.

Per ogni compilazione il manifest conserva dimensione, profondità, conteggio completo
delle operazioni, base, livello e seed di transpiler, versione delle librerie e SHA-256
della sequenza compilata. Modelli e risultati dichiarano inoltre schema, revisione,
preset, seed e definizione dell'etichetta; l'inferenza rifiuta un classificatore con
manifest incompatibile. In tal modo un artefatto non può essere associato per errore a
una revisione differente del circuito.

\subsubsection{Seed, tie-break e confronto appaiato}

Per le campagne N15, una sola simulazione fornisce l'istogramma comune a M1, TOP-4 e M2.
Il seed è separato fra repliche e iterazioni secondo
\begin{equation}
\texttt{seed\_simulator}=\texttt{rep}\cdot1\,000\,000+
\texttt{iteration}\cdot10\,000,
\end{equation}
una distanza superiore al numero di shot che evita la sovrapposizione degli stream
Aer fra semi consecutivi. La prima iterazione di successo è registrata separatamente per
ciascuna strategia; pareggi, codifica dei fallimenti e test appaiati seguono il protocollo
della Sezione~\ref{app:obiettivi:protocollo}. Il JSON schema~2 conserva input completi,
vettori appaiati, seed, test, versioni software, hash e manifest; figure e tabelle
leggono questi artefatti senza trascrizione manuale.

\subsection{Dettagli implementativi dei componenti QEC}
\label{app:dettagli_qec}

\subsubsection{M5: codice a ripetizione}

Gli stabilizzatori sono misurati indirettamente copiandone la parità su qubit ancillari,
senza misurare i dati e collassare l'informazione logica. Porte identità collocate nel
punto di iniezione costituiscono l'aggancio a cui Aer associa il canale senza modificare
lo stato ideale. Le varianti in base $Z$ e $X$ condividono lo stesso schema, con i cambi
di base necessari alla rilevazione, rispettivamente, di bit-flip e phase-flip
(Listato~\ref{lst:qec_repetition}).

\subsubsection{M6: codice di Steane}

Lo stato $\ket{0_L}$ è preparato ponendo tre qubit seme in sovrapposizione e
propagandone il supporto con CNOT. Gli stabilizzatori CSS di tipo $X$ e $Z$ alimentano
una lookup table di correzione a peso minimo. La stima di $p_L$ usa un modello di Pauli;
l'estrazione di sindrome è ideale e non costituisce un protocollo fault-tolerant
completo (Listato~\ref{lst:qec_steane}).

\subsubsection{M7: surface code}

Stim genera il circuito di memoria e campiona i detection event; PyMatching costruisce
dal detector error model il decoder \ac{MWPM}. Nell'esperimento sul crosstalk, il canale
fra dati adiacenti è inserito ricorsivamente anche nei blocchi \texttt{REPEAT}; i qubit
di misura sono identificati dal primo \texttt{MR}, non dalla misura finale che comprende
anche i dati. I Listati~\ref{lst:crosstalk} e~\ref{lst:stim_pymatching} mostrano i due
passaggi.

\subsubsection{M8/M13 e M10: proxy di Shor e decoder appreso}

Nel proxy a livelli, $p_L$ è la probabilità totale del Pauli non-identità applicato dopo
ogni gate incluso nel modello. Prima dell'esecuzione Aer viene convertito in
$\lambda_1=4p_L/3$ per i gate 1Q e $\lambda_2=16p_L/15$ per quelli 2Q; le \texttt{rz}
restano virtuali e i canali agiscono su \texttt{sx/x} e \texttt{cx}. Questo parametro
per gate non coincide con il logical error rate per round o memoria di M6/M7 senza una
compilazione fault-tolerant e un mapping esplicito. In M10, decoder appreso e
\ac{MWPM} ricevono esattamente gli stessi detection event, così il confronto riguarda la
regola di decodifica e non i campioni osservati.

\end{onehalfspace}

\subsection{Estratti dei sorgenti}

\begin{lstlisting}[caption={Installazione dell'ambiente Python e delle librerie necessarie}, label={lst:installazione_dellambiente_python_e}]
python3 -m venv ~/quantum-env
source ~/quantum-env/bin/activate
python -m pip install -r requirements.txt
python -c "import qiskit, qiskit_aer; print(qiskit.__version__)"
\end{lstlisting}

\begin{lstlisting}[caption={Implementazione della Quantum Fourier Transform inversa}, label={lst:implementazione_della_quantum_fourier}]
def inverse_qft(n_qubits):
    qc = QuantumCircuit(n_qubits, name='QFT_dagger')
    for q in range(n_qubits // 2):
        qc.swap(q, n_qubits - q - 1)
    for j in range(n_qubits):
        for m in range(j):
            qc.cp(-np.pi / 2 ** (j - m), m, j)
        qc.h(j)
    return qc
\end{lstlisting}

\begin{lstlisting}[caption={Moltiplicazione modulare N15 corretta: l'operazione $\times a$ viene ripetuta \texttt{power} volte}, label={lst:modular_exponentiation_per_n15}]
def c_amod15(a, power):
    if a not in [2, 7, 8, 11, 13]:
        raise ValueError("base non supportata per N=15")
    U = QuantumCircuit(4)
    for _ in range(power):
        if a in [2, 13]:
            U.swap(0, 1); U.swap(1, 2); U.swap(2, 3)
        if a in [7, 8]:
            U.swap(2, 3); U.swap(1, 2); U.swap(0, 1)
        if a == 11:
            U.swap(1, 3); U.swap(0, 2)
        if a in [7, 11, 13]:
            U.x(range(4))
    return U.control(annotated=False)
\end{lstlisting}

\begin{lstlisting}[caption={Circuito N15 e post-processing classico}, label={lst:circuito_completo_di_shor}]
def shor_circuit(N, a, n_count):
    if N != 15:
        return shor_circuit_beauregard(N, a, n_count)
    n_work = 4
    qc = QuantumCircuit(n_count + n_work, n_count)
    qc.h(range(n_count))
    qc.x(n_count + 3)   # |1> nella convenzione del registro textbook
    for j in range(n_count):
        power = 2 ** j
        if power % 4 != 0:       # U^4 = I per a=7 mod 15
            qc.append(c_amod15(a, power),
                      [j] + list(range(n_count, n_count + n_work)))
    qc.barrier()
    qc.append(inverse_qft(n_count), range(n_count))
    qc.measure(range(n_count), range(n_count))
    return qc

def extract_factors(measured_value, n_count, N, a):
    if measured_value == 0:
        return None, None
    phase = measured_value / 2 ** n_count
    r = Fraction(phase).limit_denominator(N).denominator
    if r % 2:
        return None, None
    for candidate in (gcd(a ** (r // 2) - 1, N),
                      gcd(a ** (r // 2) + 1, N)):
        if 1 < candidate < N:
            return candidate, N // candidate
    return None, None
\end{lstlisting}

\begin{lstlisting}[caption={Noise model uniforme illustrativo: RZ virtuale, durate 1Q/2Q separate}, label={lst:costruzione_del_noise_model}]
BASIS_GATES = ('rz', 'sx', 'x', 'cx')
GATE_TIME_1Q_NS = 50.0
GATE_TIME_2Q_NS = 300.0

def compose_errors(errors):
    out = None
    for error in errors:
        if error is not None:
            out = error if out is None else out.compose(error)
    return out

def build_noise_model(eps_1q, eps_2q, t1_ns, t2_ns, p_ro=0.02):
    # eps_1q ed eps_2q sono i lambda di depolarizing_error.
    thermal_1q = thermal_relaxation_error(
        t1_ns, t2_ns, GATE_TIME_1Q_NS)
    thermal_one_2q = thermal_relaxation_error(
        t1_ns, t2_ns, GATE_TIME_2Q_NS)
    thermal_2q = thermal_one_2q.tensor(thermal_one_2q)
    error_1q = compose_errors([
        depolarizing_error(eps_1q, 1) if eps_1q else None,
        thermal_1q,
    ])
    error_2q = compose_errors([
        depolarizing_error(eps_2q, 2) if eps_2q else None,
        thermal_2q,
    ])
    nm = NoiseModel()
    nm.add_all_qubit_quantum_error(error_1q, ['sx', 'x'])
    nm.add_all_qubit_quantum_error(error_2q, ['cx'])
    if p_ro:
        nm.add_all_qubit_readout_error(
            ReadoutError([[1-p_ro, p_ro], [p_ro, 1-p_ro]]))
    return nm
\end{lstlisting}

\begin{lstlisting}[caption={Compilazione deterministica, tie-break SHA-256 e baseline TOP-1}, label={lst:pipeline_del_metodo_1}]
def compile_shor_circuit(N, a, n_count):
    return transpile(
        shor_circuit(N, a, n_count),
        basis_gates=['rz', 'sx', 'x', 'cx'],
        optimization_level=2,
        seed_transpiler=20260819,
    )

def rank_measurements(counts, tie_seed):
    def priority(bitstring):
        payload = f'{int(tie_seed)}:{bitstring}'.encode('ascii')
        return hashlib.sha256(payload).digest()
    return sorted(counts.items(),
                  key=lambda item: (-item[1], priority(item[0])))

def method1_from_counts(counts, tie_seed, n_count, N, a):
    bitstring = rank_measurements(counts, tie_seed)[0][0]
    return extract_factors(int(bitstring, 2), n_count, N, a)
\end{lstlisting}

\begin{lstlisting}[caption={Dataset condiviso per l'audit delle etichette TOP-1 e TOP-16}, label={lst:generazione_del_dataset_con}]
LABEL_TOP_KS = (1, 16)
labels = {top_k: [] for top_k in LABEL_TOP_KS}

for sample_index in range(n_samples):
    sample_seed = seed * 1_000_000 + sample_index * 10_000
    counts = simulator.run(
        compiled, noise_model=sample_noise, shots=shots,
        seed_simulator=sample_seed,
    ).result().get_counts()
    ranked = rank_measurements(counts, tie_seed=sample_seed)
    for top_k in LABEL_TOP_KS:
        useful = any(
            extract_factors(int(bits, 2), n_count, N, a)[0] is not None
            for bits, _ in ranked[:top_k]
        )
        labels[top_k].append(int(useful))
    X.append(normalized_histogram(counts, n_count, shots))

# Ogni modello salva schema_version=2.0, label_top_k e manifest.
\end{lstlisting}

\begin{lstlisting}[caption={Decisione M2 sullo stesso istogramma usato da M1 e TOP-4}, label={lst:loop_di_inferenza_del}]
def evaluate_same_histogram(counts, simulation_seed, classifier,
                            n_count, N, a, shots):
    ranked = rank_measurements(counts, simulation_seed)
    valid = [
        extract_factors(int(bits, 2), n_count, N, a)[0] is not None
        for bits, _ in ranked[:4]
    ]
    m1_success = valid[0]
    top4_success = any(valid)
    feature = normalized_histogram(counts, n_count, shots)
    m2_success = bool(classifier.predict([feature])[0]) and top4_success
    return m1_success, top4_success, m2_success
\end{lstlisting}

\begin{lstlisting}[caption={Orchestrazione appaiata delle sole campagne rumorose N15}, label={lst:schema_di_orchestrazione_degli}]
PRESETS = {
    'reference': dict(N=15, a=7, n_count=8,
                      eps_1q=1e-3, eps_2q=1e-2,
                      t1_ns=100_000, t2_ns=80_000, p_ro=0.02),
    'stress':    dict(N=15, a=7, n_count=8,
                      eps_1q=5e-3, eps_2q=5e-2,
                      t1_ns=50_000, t2_ns=30_000, p_ro=0.05),
}
K_REPETITIONS, SHOTS, MAX_ITER = 30, 1024, 50

for rep in range(K_REPETITIONS):
    first = {'m1': None, 'top4': None, 'm2': None}
    for iteration in range(1, MAX_ITER + 1):
        simulation_seed = rep * 1_000_000 + iteration * 10_000
        counts = simulator.run(compiled, shots=SHOTS,
             seed_simulator=simulation_seed).result().get_counts()
        flags = evaluate_same_histogram(counts, simulation_seed, ...)
        update_first_success(first, flags, iteration)
    encoded = {name: value if value is not None else MAX_ITER + 1
               for name, value in first.items()}

# Inferenza: wilcoxon(paired_a, paired_b,
#                     zero_method='pratt', alternative=...)
# Output: schema_version='2.0', input espliciti e manifest/hash.
\end{lstlisting}

\begin{lstlisting}[caption={Codifica ed estrazione di sindrome del codice a ripetizione.}, label={lst:qec_repetition}]
if logical_value == 1:
    qc.x(0)
qc.cx(0, 1); qc.cx(0, 2)
if basis == 'X':
    qc.h(DATA)
for q in DATA:
    qc.id(q)                       # punto di iniezione del rumore
if basis == 'Z':
    qc.cx(0, 3); qc.cx(1, 3)
    qc.cx(1, 4); qc.cx(2, 4)
else:
    qc.h(3); qc.cx(3, 0); qc.cx(3, 1); qc.h(3)
    qc.h(4); qc.cx(4, 1); qc.cx(4, 2); qc.h(4)
qc.measure(ANC[0], 0); qc.measure(ANC[1], 1)
\end{lstlisting}

\begin{lstlisting}[caption={Preparazione di $\ket{0_L}$ e misura degli stabilizzatori di Steane.}, label={lst:qec_steane}]
def encode_zero_L(qc):
    for seed in SEED:
        qc.h(seed)
    for seed, support in SEED.items():
        for data in support:
            if data != seed:
                qc.cx(seed, data)

def measure_syndrome(qc):
    for k, support in enumerate(GZ):
        for data in support:
            qc.cx(data, ANC_Z[k])
    for k, support in enumerate(GX):
        qc.h(ANC_X[k])
        for data in support:
            qc.cx(ANC_X[k], data)
        qc.h(ANC_X[k])
\end{lstlisting}

\begin{lstlisting}[caption={Iniezione ricorsiva del crosstalk nel circuito del surface code.}, label={lst:crosstalk}]
def inject(block, pairs, p_ct):
    out = stim.Circuit()
    first = True
    for inst in block:
        if isinstance(inst, stim.CircuitRepeatBlock):
            out += inject(inst.body_copy(), pairs, p_ct) * inst.repeat_count
            continue
        out.append(inst)
        if inst.name == 'DEPOLARIZE1':
            if first:
                out.append('DEPOLARIZE2', pairs, p_ct)
            first = not first
    return out
\end{lstlisting}

\begin{lstlisting}[caption={Stima di $p_L$ con Stim e PyMatching.}, label={lst:stim_pymatching}]
circuit = stim.Circuit.generated(
    'surface_code:rotated_memory_x', distance=d, rounds=rounds,
    after_clifford_depolarization=p,
    after_reset_flip_probability=p,
    before_measure_flip_probability=p,
)
sampler = circuit.compile_detector_sampler()
events, observables = sampler.sample(shots, separate_observables=True)
model = circuit.detector_error_model(decompose_errors=True)
matching = pymatching.Matching.from_detector_error_model(model)
predictions = matching.decode_batch(events)
p_L = (predictions != observables).any(axis=1).mean()
\end{lstlisting}
